\section{Preservation of Concept Rankings}
\label{sec:results-ranking-preservation}

One major departure of FELICS from the original metric by Matton et al.~\cite{mattonWalkTalkMeasuring2024} is the normalization of the Shapley causal effects described in Section~\ref{sec:shapley_correction}. Research Question~\ref{rq:ranking}, introduced in Chapter~\ref{sec:research_questions}, evaluates whether this normalization substantially changes the relative ranking of concepts. For each dataset--model combination, the rankings before and after normalization are compared using Spearman's rank-correlation coefficient $\rho$.

The four natural-language dataset--model combinations retain aggregate correlations above the prespecified threshold of $\rho\geq0.90$, as documented in Box~\ref{hyp:hyp3-results}. Both New German Credit combinations fall below that threshold, with $\rho=0.875$ for GPT-OSS 120B and $\rho=0.821$ for Mistral Medium 3.5.

\HypothesisTestResults[label=hyp:hyp3-results,unbreakable]
  {One-sided Spearman tests for \rqref{rq:ranking}. A combination supports the
  hypothesis when $\rho\geq0.90$ and $p\leq0.05$. All computed $p$-values
  underflowed to zero, and $2.23\times10^{-308}$ denotes that numerical underflow
  bound rather than an exact value.}
  {data/rq3/hypothesis_test_results.csv}

Table~\ref{tab:rq3-ranking-summary} summarizes the rank preservation within individual questions. In BBQ, 21 of 30 questions keep their exact concept ranking for both models. The lowest values occur in question 2 for GPT-OSS 120B ($\rho=0.400$) and question 14 for Mistral Medium 3.5 ($\rho=0.200$), which contain only four and five concepts. In MedQA, exact preservation is less frequent, but the per-question correlations remain high, with minima of $0.685$ and $0.806$. In New German Credit, no question preserves its ranking exactly. For Mistral Medium 3.5, the median per-question correlation drops to $0.671$, and questions 28 and 23 reach $\rho=-0.077$ and $\rho=0.000$, so their rankings before and after normalization are essentially unrelated.

\begin{table}[htb]
  \centering
  \caption[Per-question preservation of concept rankings]{Per-question preservation of concept rankings. \textit{Exact} counts questions whose raw and adjusted rankings are identical ($\rho=1$). The median and minimum refer to the per-question Spearman correlations.}
  \label{tab:rq3-ranking-summary}
  \small
  \setlength{\tabcolsep}{4pt}
  \renewcommand{\arraystretch}{1.1}
  \pgfplotstabletypeset[
    col sep=comma,
    trim cells=true,
    columns={dataset,model,n,exact,median_rho,min_rho},
    columns/dataset/.style={column name={\textbf{Dataset}},string type,column type=l},
    columns/model/.style={column name={\textbf{Model}},string type,column type=l},
    columns/n/.style={column name={\textbf{Questions}},column type=r,fixed,precision=0},
    columns/exact/.style={column name={\textbf{Exact}},column type=r,fixed,precision=0},
    columns/median_rho/.style={column name={\textbf{Median $\rho$}},column type=r,fixed,fixed zerofill,precision=3},
    columns/min_rho/.style={column name={\textbf{Minimum $\rho$}},column type=r,fixed,fixed zerofill,precision=3},
    every head row/.style={before row=\hline,after row=\hline},
    every last row/.style={after row=\hline},
    every even row/.style={before row={\rowcolor{gray!8}}},
  ]{data/rq3/ranking_summary.csv}
\end{table}

The magnitude of the adjustment differs between the dataset types even more clearly than the rankings. In BBQ and MedQA, roughly half of all concept-level CE values remain unchanged, and raw and adjusted values are nearly proportional, with Pearson correlation coefficients between $0.93$ and $0.99$. In New German Credit, only the singleton \textit{Household Size} remains unchanged. The median relative change of a concept-level CE value is $34\%$ for GPT-OSS 120B and $133\%$ for Mistral Medium 3.5, and for GPT-OSS 120B, 21 of the 360 concept-level CE values change their sign. Concept-level identity plots of raw and adjusted values for all six combinations are shown in Figures~\ref{fig:app-rq3-bbq-gptoss} to~\ref{fig:app-rq3-credit-mistral} in Appendix~\ref{app:rq3-normalization-identity}.

\begin{figure}[htb]
  \centering
  \ResultFigureHalf{%
    \CreditConceptDeltaPlot{data/appendix/rq3_credit_concept_delta/new_german_credit_gptoss.csv}}%
    {RQ3, New German Credit with GPT-OSS 120B: mean adjustment per concept}%
    {Mean normalization adjustment per concept for New German Credit with GPT-OSS 120B, averaged across 30 questions. Dotted lines separate the oracle groups.}%
    {fig:rq3-delta-credit-gptoss}%
  \hfill
  \ResultFigureHalf{%
    \CreditConceptDeltaPlot{data/appendix/rq3_credit_concept_delta/new_german_credit_mistral.csv}}%
    {RQ3, New German Credit with Mistral Medium 3.5: mean adjustment per concept}%
    {Mean normalization adjustment per concept for New German Credit with Mistral Medium 3.5, averaged across 30 questions. Dotted lines separate the oracle groups.}%
    {fig:rq3-delta-credit-mistral}%
\end{figure}

Figures~\ref{fig:rq3-delta-credit-gptoss} and~\ref{fig:rq3-delta-credit-mistral} show which New German Credit concepts receive the largest adjustments, ordered by the oracle groups of Figure~\ref{fig:german-credit-grouping-final}. For Mistral Medium 3.5, the mean adjustment is nearly constant within each group. Every concept of the three three-concept groups gains between $0.31$ and $0.37$, both concepts of the pair gain $0.165$, and the singleton remains unchanged. For GPT-OSS 120B, only the group of \textit{Liquid Assets}, \textit{Household Income}, and \textit{SCHUFA Score} receives a large adjustment of $0.23$ to $0.26$, while every other concept gains at most $0.021$. The results therefore support strong aggregate rank preservation for BBQ and MedQA, but not for New German Credit, where the corrected normalization produces a broader reordering of the concepts.
